% Chapter 2 -- Related Work (rewritten 25 Nov 2026 for the final scope, Amendment A2)
\chapter{Related Work}
\label{ch:sota}
\lhead{Chapter 2. \emph{Related Work}}

Four lines of work meet in this thesis: evaluation of language models against regulation, the use of language models as judges, the measurement of human agreement, and the statistics of deciding with a guaranteed error, either by abstaining (selective prediction) or by stopping early (sequential inference). A fifth line, fairness audits of language-model decisions in credit and insurance, supplies the application. The chapter reviews each and ends with the gaps that define RQ1--RQ3.

\section{Regulation-facing evaluation and audits}
HELM \citep{helm} established multi-dimensional evaluation of language models; TrustLLM \citep{trustllm} and DecodingTrust \citep{decodingtrust} extend it to trust dimensions. None is organised by legal obligation. COMPL-AI \citep{complai} is the nearest neighbour: it gives a technical interpretation of the AI Act, maps it by hand onto existing benchmarks and evaluates twelve models, finding all weakest on robustness, diversity and fairness. Its mapping is manual, its suite frozen, its scoring never validated against humans, and its authors state that it is not an official audit. AIR-Bench 2024 \citep{airbench} derives 314 risk categories from eight government regulations and sixteen company policies and tests 5{,}694 human-audited prompts; it covers content safety rather than the obligations of a high-risk deployment, with no thresholds fixed before results. Bench-2-CoP \citep{bench2cop} maps 194{,}955 benchmark questions onto the Act's taxonomy and finds coverage concentrated on hallucination and reliability, with no verdicts and Annex~III out of scope. AIReg-Bench \citep{airegbench} asks LLMs to judge technical documentation for AI Act violations against legal experts' labels; it judges documents, not the behaviour of a deployed model. On the audit side, SMACTR \citep{smactr} and the three-layered approach of M\"okander et al.\ \citep{mokander} organise audits around artefacts and governance layers and call for behavioural testing without supplying a method for it.

\emph{What is missing:} requirements derived from the legal text and signed before data, behavioural verdicts with stated error, and evidence that a procedure carries from one Annex~III use case to another.

\section{Language models as judges}
Zheng et al.\ \citep{zheng} introduced LLM-as-a-judge with MT-Bench and Chatbot Arena and found that GPT-4 agrees with humans about as often as humans agree with each other, while documenting position, verbosity and self-enhancement bias. The biases have since been measured directly: the same answer pair receives different verdicts depending on its order \citep{wang2023fair}; across 15 judges, 22 tasks and more than 150{,}000 evaluations position consistency stays below 0.75 even for strong judges \citep{shi2024position}; longer answers are preferred regardless of quality \citep{saito2023}; judges favour text that resembles their own \citep{wataoka2024} and models related to them by training data or family (preference leakage, \citealp{li2025leakage}). Reliability is a separate problem: deterministic decoding does not make verdicts reliable, so a single run is a weak basis for measurement \citep{schroeder2024}. On hard objective-correctness pairs strong judges barely beat chance \citep{judgebench}; judges can be swayed by manipulated answers \citep{thakur2024}; and in a legal task, models' reasons diverge from legal experts' despite good faithfulness scores \citep{namazov2026}. Surveys \citep{gusurvey} and a recent systematic review \citep{shrestha2026} converge on one conclusion: biases are systematic, single-pass evaluation is insufficient, and there is no standard way to check a judge before it is trusted.

\emph{What is missing:} a judge whose admission is decided by a measured human ceiling on a legal criterion, and whose individual verdicts carry an error guarantee.

\section{Human agreement and the validation of judges}
Agreement between raters is measured with chance-corrected coefficients: Cohen's $\kappa$ for two raters \citep{cohen1960}, Fleiss' $\kappa$ for many \citep{fleiss1971}, read against the bands of Landis and Koch \citep{landiskoch}, in which 0.61--0.80 is ``substantial''. When one label dominates, high raw agreement coexists with low $\kappa$ (the $\kappa$ paradox, \citealp{feinstein}); Gwet's AC1 \citep{gwet} corrects for chance in a way that is less sensitive to prevalence. Artstein and Poesio \citep{artstein2008} argue that agreement must be established before annotations are used as a reference, and Shankar et al.\ \citep{shankar2024} show that graders form their criteria while grading (criteria drift), which is a reason to freeze a decidable guideline before rating.

Recent work validates judges against humans but takes the human labels as the truth. The Alternative Annotator Test \citep{calderon2025} asks whether an LLM can replace an annotator using at least three annotators on a subset; JUDGE-BENCH \citep{bavaresco2025} finds that judge reliability depends on the property judged, the annotators' expertise and the text source; across about 541{,}000 judgments chance-corrected agreement sits 33--41 points below raw agreement \citep{norman2026}; and in expert domains judge--expert agreement is 60--68\%, below the experts' agreement with each other \citep{szymanski2025}. The GRAIL case study adds a direct counter-example to the assumption: four raters agreed at chance level (Fleiss $\kappa=0.07$), largely because of prevalence and an undecidable guideline (Section~\ref{sec:annotation}).

\emph{What is missing:} validation that measures the human ceiling first, for the specific legal criterion, and judges the judge relative to it.

\section{Selective prediction and certified automation}
Selective prediction lets a system abstain. Chow \citep{chow1970} characterised the optimal trade-off between error and rejection; El-Yaniv and Wiener \citep{elyaniv2010} developed the theory of selective classification; Geifman and El-Yaniv \citep{geifman} calibrated a threshold on a confidence score so that the risk among accepted items is bounded, with an exact binomial bound \citep{clopper1934}. Distribution-free methods generalise this: conformal prediction \citep{angelopoulos} and Learn then Test \citep{learntest}, which chooses a threshold by multiple testing so that a risk is controlled with high probability, given exchangeable data. For LLM judges, Trust or Escalate \citep{jung2025} gives a provable guarantee of human agreement on the judgments it accepts and escalates the rest to stronger judges; SCOPE \citep{badshah2026} calibrates an acceptance threshold for pairwise judgments under exchangeability. Prediction-powered inference \citep{ppi} and confidence-driven inference \citep{gligoric2025} use model labels to reduce the number of human labels for a valid population estimate, and Dorner et al.\ \citep{dorner} bound how much a judge no better than the evaluated model can save.

\emph{What is missing:} the guarantee applied to verdicts derived from a legal criterion, after the reference itself has been shown reliable, with the assumptions (sampling, frozen judge, threshold selection, reference error) stated as part of the claim.

\section{Sequential inference}
Sequential tests decide while data arrive. Wald's sequential probability ratio test \citep{wald1945} is the origin; group-sequential designs in clinical trials control the type-I error across planned interim looks with constant boundaries \citep{pocock1977} or boundaries that are strict early and relaxed late \citep{obrien1979}, and $\alpha$-spending functions make the number and timing of looks flexible \citep{landemets}. Confidence sequences are valid at any stopping time \citep{howard}, and betting-based constructions give tight anytime-valid bounds for bounded data \citep{waudbysmith2024}. In machine-learning auditing, Chugg et al.\ \citep{chugg2023} audit the fairness of deployed models by betting, Richter et al.\ \citep{richter2025} detect behaviour change between model versions, and recent work evaluates LLMs anytime-validly with fewer samples \citep{celeus2026,hsu2026}.

\emph{What is missing:} each of these controls one metric, one score or one change at a time. An audit shares one budget across many requirements and must return three-way verdicts (PASS, FAIL, UNDETERMINED); whether early stopping across them preserves the verdicts of the full audit is untested.

\section{Fairness of LLM decisions in credit and insurance}
Paired designs hold an applicant fixed and vary one protected signal. In mortgage underwriting, LLMs deny Black applicants more often and price them higher, most strongly at low credit scores \citep{bowen}; race-based discrepancies appear on counterfactual applicants identical except for race \citep{cook2025}; and fair outputs can hide race information that persists inside the model and reverses decisions when re-injected \citep{tripathy2026}, a limit of any output-only audit. McNemar's exact test \citep{mcnemar} measures differential treatment of the same applicant, while the four-fifths rule \citep{eeoc} compares group rates; the case study showed that the two disagree completely on the same answers. In insurance, the Court of Justice ruled in \emph{Test-Achats} \citep{testachats} that sex may not change individual premiums, from 21~December 2012. An audit of six LLMs on 1{,}388 motor claims finds no male/female gap under full information but gaps under other prompt set-ups \citep{huang2026binary}; conventional life and health pricing models show income-related premium gaps against a 10\% tolerance \citep{gupta2025}; and discrimination-free pricing constructs prices that do not depend on the protected attribute \citep{lindholm2022}. Equivalence testing with two one-sided tests \citep{schuirmann1987} supplies the statistical form for showing that a gap is small rather than merely not significant.

\emph{What is missing:} an individual-level audit of LLMs that price life or health insurance, with a margin fixed in advance and real applicant profiles.

\section{Gaps, research questions and calibration}
\begin{table}[H]
\centering
\caption{Gaps in the literature and the research question or calibration step that answers each.}
\label{tab:gaps}
\small
\begin{tabularx}{\textwidth}{@{}lXl@{}}
\toprule
Gap & What is missing & Answered by \\
\midrule
G1 & Judges are validated against human labels that are never shown to be reliable & CAL1, CAL2a \\
G2 & Guaranteed-error judges exist only for general benchmarks, never for a legal criterion with stated assumptions & CAL2b \\
G3 & Valid early stopping exists for one metric at a time, not for many audit requirements with three-way verdicts & RQ1 \\
G4 & Legal text is mapped onto benchmarks by hand; nothing shows a procedure carries to a second Annex~III use case & RQ2 (H2) \\
G5 & No individual-level audit of LLMs that price life or health insurance & RQ3 (H3a--H3c) \\
\bottomrule
\end{tabularx}
\end{table}

\section*{Literature survey in tabular form}
\addcontentsline{toc}{section}{Literature survey in tabular form}
\par\vspace{0.5em}
{\scriptsize
\setlength{\tabcolsep}{3pt}\setlength{\LTpre}{1.5em}\setlength{\LTpost}{1em}
\renewcommand{\arraystretch}{1.15}
\begin{longtable}{@{}>{\raggedright\arraybackslash}p{2.6cm}>{\raggedright\arraybackslash}p{1.8cm}>{\raggedright\arraybackslash}p{4.3cm}>{\raggedright\arraybackslash}p{4.3cm}>{\raggedright\arraybackslash}p{1.3cm}@{}}
\caption{Literature survey: what each work shows, what it leaves open, and the research question or calibration step it informs.}\label{tab:litsurvey}\\
\toprule
\textbf{Work} & \textbf{Area} & \textbf{What it shows} & \textbf{What it leaves open for GRAIL} & \textbf{RQ/CAL} \\
\midrule
\endfirsthead
\toprule
\textbf{Work} & \textbf{Area} & \textbf{What it shows} & \textbf{What it leaves open for GRAIL} & \textbf{RQ/CAL} \\
\midrule
\endhead
\bottomrule
\endfoot
COMPL-AI \citep{complai} & Regulation & AI Act mapped by hand onto benchmarks; no model fully compliant & Manual mapping, no validated scoring, no verdict error & RQ2 \\
AIR-Bench 2024 \citep{airbench} & Regulation & 314 risk categories from regulations and policies & Content safety, not high-risk obligations & RQ2 \\
Bench-2-CoP \citep{bench2cop} & Regulation & Benchmark coverage of the Act's taxonomy & Coverage, no verdicts; Annex III out of scope & RQ2 \\
AIReg-Bench \citep{airegbench} & Regulation & LLMs judging documentation against legal experts & Documents, not behaviour & RQ2 \\
\citet{zheng} & Judges & Judge--human agreement near human--human; three biases & Measured once, globally & CAL2 \\
\citet{wang2023fair}; \citet{shi2024position} & Judges & Order changes verdicts; position consistency $<0.75$ & Pairwise; no guarantee & opt. \\
\citet{saito2023}; \citet{wataoka2024}; \citet{li2025leakage} & Judges & Verbosity, self-preference and preference-leakage biases & Detected, not controlled & opt. \\
\citet{schroeder2024} & Judges & One run at temperature 0 is not reliable & No decision rule & CAL2 \\
JudgeBench \citep{judgebench} & Judges & Strong judges near chance on hard pairs & General tasks & CAL2 \\
\citet{shrestha2026} & Judges (review) & No standard pre-deployment check of a judge & A map, not evidence & CAL2 \\
\citet{feinstein}; \citet{gwet} & Agreement & $\kappa$ paradox; AC1 & --- & CAL1 \\
\citet{artstein2008}; \citet{shankar2024} & Agreement & Establish agreement first; criteria drift & How to freeze a legal criterion & CAL1 \\
\citet{calderon2025}; \citet{bavaresco2025} & Validation & Tests of replacing annotators; reliability varies by task & Human labels assumed correct & CAL1 \\
\citet{norman2026}; \citet{szymanski2025} & Validation & Chance-corrected agreement far below raw; expert domains weaker & Reference not measured first & CAL1 \\
\citet{geifman}; \citet{chow1970} & Selective & Bounded risk on accepted items; error--reject trade-off & Classifiers, i.i.d. data & CAL2 \\
Learn then Test \citep{learntest} & Selective & Threshold chosen by multiple testing for risk control & Exchangeability assumed & CAL2 \\
Trust or Escalate \citep{jung2025}; SCOPE \citep{badshah2026} & Selective & Guaranteed agreement / error for accepted judgments & Labels assumed correct; general tasks & CAL2 \\
\citet{ppi}; \citet{dorner} & Selective & Fewer labels for valid estimates; limits of judges & Rates, not individual verdicts & CAL2 \\
\citet{obrien1979}; \citet{landemets} & Sequential & Group-sequential boundaries, $\alpha$-spending & Single trials & RQ1 \\
\citet{howard}; \citet{waudbysmith2024} & Sequential & Anytime-valid confidence sequences & Conservative at audit sizes & RQ1 \\
\citet{chugg2023}; \citet{richter2025} & Sequential & Sequential fairness audit; behaviour-shift test & One metric; tests for unfairness or change & RQ1 \\
CELEUS \citep{celeus2026}; \citet{hsu2026} & Sequential & Anytime-valid LLM evaluation with fewer samples & One score per model & RQ1 \\
\citet{bowen}; \citet{cook2025} & Credit & Paired audits find racial disparities in LLM lending & No clause trace, no verdict reliability & RQ3 \\
\citet{tripathy2026} & Credit & Fair outputs, biased internals & Output-only audits miss latent bias & limit \\
\citet{huang2026binary} & Insurance & Gender gaps in LLM claim decisions depend on set-up & Motor claims, no thresholds & RQ3 \\
\citet{gupta2025}; \citet{lindholm2022} & Insurance & Group pricing gaps; discrimination-free prices & Conventional models, group metrics & RQ3 \\
\emph{Test-Achats} \citep{testachats} & Law & Sex may not change individual premiums (2012) & No testing method & RQ3 \\
\end{longtable}
}
